%% Standard Library
\chapter{Standard Library}
\label{ch:stdlib}

The Lambdust standard library extends R7RS-small with additional modules for advanced programming while maintaining complete backward compatibility. All R7RS-small procedures are available unchanged.

\section{Library Organization}
\label{sec:library-organization}

The standard library is organized into modules:

\begin{itemize}
\item \textbf{R7RS Core}: Complete R7RS-small implementation
\item \textbf{Extended Base}: Enhanced versions of core libraries
\item \textbf{Gradual Types}: Type system integration
\item \textbf{Effects}: Algebraic effects and handlers
\item \textbf{Concurrency}: Actor model and parallelism
\item \textbf{Collections}: Advanced data structures
\item \textbf{Text}: Unicode and pattern matching
\item \textbf{I/O}: Enhanced input/output operations
\item \textbf{FFI}: Foreign function interface
\item \textbf{Meta}: Reflection and metaprogramming
\end{itemize}

\section{R7RS-small Compatibility}
\label{sec:r7rs-compatibility}

\indexentry{R7RS compatibility}
\indexentry{base library}

All R7RS-small libraries are fully supported:

\begin{example}[R7RS Libraries]
\begin{lambdustcode}
;; Import R7RS libraries
(import (scheme base)
        (scheme case-lambda)
        (scheme char)
        (scheme complex)
        (scheme cxr)
        (scheme eval)
        (scheme file)
        (scheme inexact)
        (scheme lazy)
        (scheme load)
        (scheme process-context)
        (scheme read)
        (scheme repl)
        (scheme time)
        (scheme write))

;; All R7RS procedures work unchanged
(define (r7rs-example)
  (let* ((numbers '(3 1 4 1 5 9 2 6))
         (doubled (map (lambda (x) (* x 2)) numbers))
         (filtered (filter (lambda (x) (> x 5)) doubled)))
    (fold-left + 0 filtered)))

;; R7RS file operations
(call-with-output-file "output.txt"
  (lambda (port)
    (write '(hello world) port)
    (newline port)))

;; R7RS parameter objects
(parameterize ((current-output-port (open-output-string)))
  (display "Hello")
  (get-output-string (current-output-port)))
\end{lambdustcode}
\end{example}

\section{Extended Base Library}
\label{sec:extended-base}

\indexentry{extended base}

The \library{lambdust base} library extends R7RS base:

\subsection{Enhanced List Operations}
\label{subsec:enhanced-lists}

\begin{procspec}{(take lst k)}
Returns the first \var{k} elements of \var{lst}.
\end{procspec}

\begin{procspec}{(drop lst k)}
Returns \var{lst} with the first \var{k} elements removed.
\end{procspec}

\begin{procspec}{(split-at lst k)}
Returns two values: \code{(take lst k)} and \code{(drop lst k)}.
\end{procspec}

\begin{example}[Enhanced List Operations]
\begin{lambdustcode}
(import (lambdust base))

;; List utilities
(take '(1 2 3 4 5) 3)        ; → (1 2 3)
(drop '(1 2 3 4 5) 2)        ; → (3 4 5)
(split-at '(1 2 3 4 5) 2)    ; → (1 2) (3 4 5)

;; List comprehensions
(list-comprehension (x * x) (x in '(1 2 3 4 5)) (x > 2))
; → (9 16 25)

;; Lazy lists
(define fibs
  (lazy-list 0 1 (lazy-map + fibs (lazy-tail fibs))))

(take (lazy-list->list fibs) 10)
; → (0 1 1 2 3 5 8 13 21 34)

;; List zipping
(zip '(1 2 3) '(a b c) '(x y z))     ; → ((1 a x) (2 b y) (3 c z))
(zip-with + '(1 2 3) '(4 5 6))      ; → (5 7 9)

;; Grouping and partitioning  
(group-by even? '(1 2 3 4 5 6))     ; → ((#f 1 3 5) (#t 2 4 6))
(partition even? '(1 2 3 4 5 6))    ; → (2 4 6) (1 3 5)
\end{lambdustcode}
\end{example}

\subsection{Enhanced String Operations}
\label{subsec:enhanced-strings}

\begin{procspec}{(string-split str delimiter)}
Splits \var{str} on \var{delimiter} returning a list of substrings.
\end{procspec}

\begin{procspec}{(string-join strings separator)}
Joins \var{strings} with \var{separator}.
\end{procspec}

\begin{example}[Enhanced String Operations]
\begin{lambdustcode}
;; String splitting and joining
(string-split "hello,world,scheme" ",")  ; → ("hello" "world" "scheme")
(string-join '("hello" "world") " ")     ; → "hello world"

;; String formatting
(format "Hello, ~a! You have ~d messages." "Alice" 5)
; → "Hello, Alice! You have 5 messages."

;; String templates
(string-template "Hello, {name}!" '((name . "Bob")))
; → "Hello, Bob!"

;; Case conversion
(string-titlecase "hello world")         ; → "Hello World"
(string-capitalize "hello")              ; → "Hello"

;; String predicates
(string-prefix? "hello" "hello world")   ; → #t
(string-suffix? "world" "hello world")   ; → #t
(string-contains? "ell" "hello")         ; → 1

;; Unicode operations
(string-normalize "café" 'nfc)           ; NFC normalization
(string-grapheme-count "👨‍👩‍👧‍👦")           ; → 1 (family emoji)
\end{lambdustcode}
\end{example}

\section{Type System Integration}
\label{sec:type-integration}

\indexentry{type library}

The \library{lambdust types} library provides type system integration:

\begin{example}[Type System Library]
\begin{lambdustcode}
(import (lambdust types))

;; Type predicates
(define-predicate natural? (lambda (x) (and (integer? x) (>= x 0))))
(define-predicate positive? (lambda (x) (and (real? x) (> x 0))))

;; Type assertions
(assert-type 42 Natural)                 ; Succeeds
(assert-type -5 Natural)                 ; Error

;; Type conversion
(cast-type "123" String Number)          ; → 123
(safe-cast "123" Number)                 ; → (Just 123)
(safe-cast "abc" Number)                 ; → Nothing

;; Contract checking
(define/contract factorial
  (-> Natural Natural)
  (lambda (n)
    (if (<= n 1) 1 (* n (factorial (- n 1))))))

;; Higher-order contracts
(define/contract map-numbers
  (-> (-> Number Number) (List Number) (List Number))
  (lambda (f lst) (map f lst)))

;; Refinement types  
(define NonEmptyList (Refinement (List Any) (lambda (lst) (not (null? lst)))))
(define/contract safe-car
  (-> NonEmptyList Any)
  (lambda (lst) (car lst)))

;; Gradual typing utilities
(dynamic->typed 42)                      ; Dynamic → typed value
(typed->dynamic (the Natural 42))        ; Typed → dynamic value
\end{lambdustcode}
\end{example}

\section{Effects Library}
\label{sec:effects-library}

\indexentry{effects library}

The \library{lambdust effects} library provides effect operations:

\begin{example}[Effects Library]
\begin{lambdustcode}
(import (lambdust effects))

;; Built-in effects
(define (stateful-computation)
  (perform State set 0)
  (perform State modify +1)
  (perform State modify +1)
  (perform State get))

;; Exception effects  
(define (safe-division x y)
  (if (= y 0)
      (perform Exception raise "Division by zero")
      (/ x y)))

;; I/O effects
(define (interactive-program)
  (perform IO write-line "Enter your name:")
  (let ((name (perform IO read-line)))
    (perform IO write-line (string-append "Hello, " name "!"))))

;; Custom effects
(define-effect Logger
  (log :: (-> String Unit))
  (set-level :: (-> Symbol Unit)))

(define (logged-computation)
  (perform Logger set-level 'debug)
  (perform Logger log "Starting computation")
  (let ((result (complex-calculation)))
    (perform Logger log "Computation complete")
    result))

;; Effect composition
(define (multi-effect-computation)
  (perform Logger log "Starting")
  (perform State set 0)
  (let ((input (perform IO read-line)))
    (let ((value (string->number input)))
      (if value
          (begin
            (perform State set value)
            (perform Logger log "Value set")
            (perform State get))
          (perform Exception raise "Invalid input")))))

;; Effect handlers
(define (run-with-state-and-logging initial computation)
  (with-state-handler initial
    (with-logging-handler 'console
      (with-exception-handler
        (lambda (e) (display "Error: ") (display e) (newline))
        computation))))
\end{lambdustcode}
\end{example}

\section{Concurrency Library}
\label{sec:concurrency-library}

\indexentry{concurrency library}

The \library{lambdust concurrent} library provides concurrency primitives:

\begin{example}[Concurrency Library]
\begin{lambdustcode}
(import (lambdust concurrent))

;; Actor operations
(define (worker-actor id)
  (let loop ()
    (receive
      (('work data reply-to)
       (let ((result (process-data data)))
         (send reply-to result))
       (loop))
      (('shutdown) 'done))))

(define worker (spawn worker-actor 1))
(send worker `(work ,some-data ,self))

;; Future operations
(define future-1 (future (expensive-computation-1)))
(define future-2 (future (expensive-computation-2)))

;; Wait for both
(let ((result-1 (await future-1))
      (result-2 (await future-2)))
  (+ result-1 result-2))

;; Parallel operations
(define results (parallel-map expensive-function data-list))

;; STM operations
(define account1 (make-tvar 1000))
(define account2 (make-tvar 500))

(define (transfer amount from to)
  (atomic
    (let ((from-bal (tvar-read from))
          (to-bal (tvar-read to)))
      (when (>= from-bal amount)
        (tvar-write! from (- from-bal amount))
        (tvar-write! to (+ to-bal amount))))))

;; Channel operations
(define numbers-channel (make-channel))
(spawn (lambda () 
         (do ((i 0 (+ i 1))) ((= i 10))
           (channel-put! numbers-channel i))))

(let loop ((sum 0) (count 0))
  (if (= count 10)
      sum
      (loop (+ sum (channel-get numbers-channel)) (+ count 1))))
\end{lambdustcode}
\end{example}

\section{Collections Library}
\label{sec:collections-library}

\indexentry{collections}
\indexentry{data structures}

The \library{lambdust collections} library provides advanced data structures:

\begin{example}[Collections Library]
\begin{lambdustcode}
(import (lambdust collections))

;; Sets
(define s1 (set 1 2 3 4))
(define s2 (set 3 4 5 6))
(set-union s1 s2)                        ; → {1 2 3 4 5 6}
(set-intersection s1 s2)                 ; → {3 4}
(set-difference s1 s2)                   ; → {1 2}

;; Maps/Dictionaries  
(define ages (map-from-alist '(("Alice" . 25) ("Bob" . 30))))
(map-ref ages "Alice")                   ; → 25
(map-set ages "Charlie" 35)              ; → new map with Charlie
(map-keys ages)                          ; → ("Alice" "Bob")

;; Bags/Multisets
(define bag1 (bag 'a 'b 'b 'c))
(bag-count bag1 'b)                      ; → 2
(bag-add bag1 'b)                        ; → bag with 3 b's

;; Sequences
(define seq (sequence 1 2 3 4 5))
(sequence-ref seq 2)                     ; → 3
(sequence-append seq (sequence 6 7))     ; → (1 2 3 4 5 6 7)

;; Persistent vectors
(define v1 (vector-from-list '(1 2 3 4 5)))
(define v2 (vector-set v1 2 'new))       ; v1 unchanged
(vector->list v2)                        ; → (1 2 new 4 5)

;; Finger trees
(define ft (finger-tree 1 2 3 4 5))
(finger-tree-split-at ft 2)              ; → (1 2) (3 4 5)

;; Priority queues
(define pq (priority-queue < 5 3 8 1 9))
(priority-queue-min pq)                  ; → 1
(priority-queue-remove-min pq)           ; → queue without 1

;; Tries
(define trie (trie-from-alist '(("cat" . 1) ("car" . 2) ("card" . 3))))
(trie-keys-with-prefix trie "ca")        ; → ("cat" "car" "card")
\end{lambdustcode}
\end{example}

\section{Text Processing}
\label{sec:text-processing}

\indexentry{text processing}
\indexentry{regular expressions}

The \library{lambdust text} library provides text processing:

\begin{example}[Text Processing]
\begin{lambdustcode}
(import (lambdust text))

;; Regular expressions
(define email-regex (regex "([\\w.-]+)@([\\w.-]+)\\.(\\w+)"))
(regex-match email-regex "alice@example.com")  
; → #<match "alice@example.com" ("alice" "example" "com")>

(regex-replace email-regex "Contact: alice@example.com" 
                "Email: $1 at $2 dot $3")
; → "Contact: Email: alice at example dot com"

(regex-find-all (regex "\\b\\w+\\b") "Hello world from Lambdust")
; → ("Hello" "world" "from" "Lambdust")

;; String patterns
(define phone-pattern (pattern "({3d})-{3d}-{4d}"))
(pattern-match phone-pattern "(555)-123-4567")
; → ((area-code . "555") (exchange . "123") (number . "4567"))

;; Text templates
(define template (text-template "Dear {{name}},\n\nYour balance is ${{balance}}."))
(template-render template 
  '((name . "Alice") (balance . "1,234.56")))

;; Parsing combinators
(define number-parser 
  (sequence (optional (char #\-))
           (one-or-more (char-range #\0 #\9))
           (optional (sequence (char #\.)
                              (one-or-more (char-range #\0 #\9))))))

(parse number-parser "-123.45")         ; → -123.45

;; CSV parsing
(csv-read-string "name,age,city\nAlice,25,Boston\nBob,30,Seattle")
; → (("name" "age" "city") ("Alice" "25" "Boston") ("Bob" "30" "Seattle"))

;; Markdown parsing
(markdown->html "# Hello\n\nThis is **bold** text.")
; → "<h1>Hello</h1>\n<p>This is <strong>bold</strong> text.</p>"
\end{lambdustcode}
\end{example}

\section{I/O Library}
\label{sec:io-library}

\indexentry{I/O library}
\indexentry{input output}

The \library{lambdust io} library provides enhanced I/O:

\begin{example}[I/O Library]
\begin{lambdustcode}
(import (lambdust io))

;; File operations
(with-input-from-file "data.txt"
  (lambda ()
    (let ((lines (read-lines)))
      (filter (lambda (line) (string-contains? line "important")) lines))))

;; Binary I/O
(with-binary-input-from-file "image.png"
  (lambda ()
    (let ((header (read-bytes 8)))
      (if (png-header? header)
          (read-png-data)
          (error "Not a PNG file")))))

;; Streaming I/O
(define (process-large-file filename processor)
  (with-input-stream (file-stream filename)
    (stream-fold processor (stream-initial-value) 
                 (stream-lines (current-input-stream)))))

;; Network I/O
(with-tcp-connection "api.example.com" 80
  (lambda (socket)
    (write-line "GET /data HTTP/1.0" socket)
    (write-line "" socket)
    (read-http-response socket)))

;; Asynchronous I/O
(define (async-file-read filename)
  (async-with-input-from-file filename
    (lambda ()
      (read-string (file-size filename)))))

;; Directory operations
(directory-list "/usr/local/bin")         ; → list of files
(directory-walk "/home/user" 
  (lambda (path type) 
    (when (eq? type 'file)
      (display path) (newline))))

;; Path manipulation
(path-join "/home" "user" "documents" "file.txt")
; → "/home/user/documents/file.txt"
(path-extension "document.pdf")          ; → "pdf"
(path-basename "/path/to/file.txt")      ; → "file.txt"

;; Temporary files
(with-temporary-file "temp-data-"
  (lambda (temp-path)
    (with-output-to-file temp-path
      (lambda () (write large-data)))
    (process-file temp-path)))
\end{lambdustcode}
\end{example}

\section{FFI Library}
\label{sec:ffi-library}

\indexentry{FFI library}

The \library{lambdust ffi} library provides foreign function interface:

\begin{example}[FFI Library] 
\begin{lambdustcode}
(import (lambdust ffi))

;; Basic foreign functions
(foreign-declare "libm.so"
  (sin :: (-> Float64 Float64))
  (cos :: (-> Float64 Float64)))

;; Structure definitions
(define-c-struct Point
  (x :: Float64)
  (y :: Float64))

;; Using foreign functions
(with-capability 'math-library
  (lambda ()
    (let ((angle 1.57))  ; π/2
      (display "sin(π/2) = ") 
      (display (foreign-call sin angle))
      (newline))))

;; Memory management
(define (process-with-buffer size)
  (let ((buffer (allocate-buffer size)))
    (dynamic-wind
      (lambda () #f)
      (lambda () (process-buffer buffer))
      (lambda () (free-buffer buffer)))))

;; Callback functions
(define compare-function
  (make-c-callback (-> Pointer Pointer Int32)
    (lambda (a b)
      (let ((val-a (pointer->integer a))
            (val-b (pointer->integer b)))
        (cond ((< val-a val-b) -1)
              ((> val-a val-b) 1)
              (else 0))))))

;; Array marshaling
(define numbers '(3 1 4 1 5 9))
(let ((c-array (list->c-array numbers Int32)))
  (with-capability 'libc
    (lambda ()
      (foreign-call qsort c-array 6 4 compare-function)))
  (c-array->list c-array))

;; Library loading
(define graphics-lib (load-dynamic-library "libgraphics.so"))
(define draw-line (get-foreign-procedure graphics-lib "draw_line"
                    (-> Float64 Float64 Float64 Float64 Void)))
\end{lambdustcode}
\end{example}

\section{Metaprogramming Library}
\label{sec:meta-library}

\indexentry{metaprogramming}
\indexentry{reflection}

The \library{lambdust meta} library provides reflection and metaprogramming:

\begin{example}[Metaprogramming Library]
\begin{lambdustcode}
(import (lambdust meta))

;; Reflection
(procedure-name +)                       ; → "+"
(procedure-arity +)                      ; → (0 . #f) ; 0 or more args
(procedure-source factorial)             ; → source code if available

;; Environment inspection
(current-environment-variables)          ; → list of bound variables
(variable-type 'x (current-environment)) ; → type of variable x
(variable-defined? 'undefined-var)       ; → #f

;; Code generation
(define-code-generator simple-loop
  (lambda (var start end body)
    `(let loop ((,var ,start))
       (if (>= ,var ,end)
           'done
           (begin ,body (loop (+ ,var 1)))))))

(simple-loop i 0 10 (display i))

;; Macro introspection  
(macro-transformer? when)                ; → #t
(expand-macro-once '(when (> x 0) (display "positive")))
; → (if (> x 0) (begin (display "positive")))

;; Runtime compilation
(define dynamic-procedure 
  (compile-lambda '(lambda (x y) (+ (* x x) (* y y)))))

(dynamic-procedure 3 4)                  ; → 25

;; Module introspection
(module-exports '(scheme base))          ; → list of exported names
(module-imports '(myapp main))           ; → list of imported modules

;; Performance profiling
(with-profiler
  (lambda ()
    (expensive-computation))
  (lambda (profile)
    (display-profile-report profile)))

;; Memory introspection
(gc-stats)                               ; → GC statistics
(object-size my-large-structure)         ; → approximate size in bytes
(heap-usage)                             ; → current heap usage
\end{lambdustcode}
\end{example}

\section{SRFI Support}
\label{sec:srfi-support}

\indexentry{SRFI}
\indexentry{Scheme Requests for Implementation}

Lambdust provides extensive SRFI support:

\begin{example}[SRFI Libraries]
\begin{lambdustcode}
;; SRFI-1: List library
(import (srfi 1))
(take '(1 2 3 4 5) 3)                   ; → (1 2 3)
(unfold negative? (lambda (x) x) (lambda (x) (- x 1)) 5)
; → (5 4 3 2 1 0)

;; SRFI-13: String library  
(import (srfi 13))
(string-tokenize "Hello, world!")       ; → ("Hello" "world")
(string-pad-right "hello" 10)           ; → "hello     "

;; SRFI-41: Streams
(import (srfi 41))
(define nats (stream-iterate (lambda (x) (+ x 1)) 0))
(stream->list (stream-take nats 10))     ; → (0 1 2 3 4 5 6 7 8 9)

;; SRFI-113: Sets and bags
(import (srfi 113))
(define s (set number-comparator 1 2 3))
(set-contains? s 2)                      ; → #t

;; SRFI-121: Generators
(import (srfi 121))
(define g (list->generator '(1 2 3 4 5)))
(generator->list (gtake g 3))            ; → (1 2 3)

;; SRFI-128: Comparators  
(import (srfi 128))
(define string-ci-comparator
  (make-comparator string? string-ci=? string-ci<? string-ci-hash))

;; SRFI-151: Bitwise operations
(import (srfi 151))
(bitwise-and 6 10)                       ; → 2
(arithmetic-shift 1 10)                  ; → 1024

;; SRFI-158: Generators and accumulators
(import (srfi 158))
(define acc (list-accumulator))
(acc 1) (acc 2) (acc 3)
(acc)                                    ; → (1 2 3)
\end{lambdustcode}
\end{example}

\section{Testing Library}
\label{sec:testing-library}

\indexentry{testing}
\indexentry{unit tests}

The \library{lambdust test} library provides testing facilities:

\begin{example}[Testing Library]
\begin{lambdustcode}
(import (lambdust test))

;; Basic assertions
(test-assert (> 5 3))
(test-equal 4 (+ 2 2))
(test-approximate 3.14159 22/7 0.01)

;; Test groups
(test-group "arithmetic tests"
  (test-equal "addition" 7 (+ 3 4))
  (test-equal "multiplication" 12 (* 3 4))
  (test-equal "division" 2 (/ 8 4)))

;; Property-based testing
(test-property "addition is commutative"
  (lambda (x y) (= (+ x y) (+ y x)))
  (generate-integers -100 100))

;; Test fixtures
(define-test-fixture database-tests
  (setup (lambda () (init-test-database)))
  (teardown (lambda () (cleanup-test-database)))
  
  (test-case "user creation"
    (let ((user (create-user "testuser")))
      (test-assert (user? user))
      (test-equal "testuser" (user-name user))))
  
  (test-case "user deletion"
    (create-user "tempuser")
    (delete-user "tempuser")
    (test-equal #f (find-user "tempuser"))))

;; Mock objects
(define-mock database-connection
  (query (lambda (sql) '((id 1) (name "Alice"))))
  (execute (lambda (sql) 'success)))

(with-mock database-connection
  (test-equal '((id 1) (name "Alice"))
              (database-query "SELECT * FROM users")))

;; Benchmark tests
(test-benchmark "list operations"
  (benchmark "append" (lambda () (append (range 0 1000) (range 1000 2000))))
  (benchmark "map" (lambda () (map (lambda (x) (* x 2)) (range 0 1000)))))

;; Coverage analysis
(with-coverage-analysis
  (run-all-tests)
  (generate-coverage-report))
\end{lambdustcode}
\end{example}

\section{Debugging Library}
\label{sec:debugging-library}

\indexentry{debugging}
\indexentry{tracing}

The \library{lambdust debug} library provides debugging tools:

\begin{example}[Debugging Library]
\begin{lambdustcode}
(import (lambdust debug))

;; Tracing
(trace factorial)
(factorial 5)
;; Trace output:
;; |(factorial 5)
;; | (factorial 4)
;; | |(factorial 3)
;; | | (factorial 2)
;; | | |(factorial 1)
;; | | |1
;; | | 1
;; | |2
;; | 6
;; |120

;; Breakpoints
(define (buggy-function x)
  (let ((intermediate (complex-calculation x)))
    (breakpoint "checking intermediate value" intermediate)
    (final-calculation intermediate)))

;; Step debugging
(step-debug
  (lambda ()
    (let ((x 10))
      (let ((y (* x 2)))
        (+ x y)))))

;; Profiling
(profile fibonacci 30)
;; Shows time spent in each function call

;; Memory debugging
(with-memory-debugging
  (lambda ()
    (let ((large-list (make-list 100000 'item)))
      (process-list large-list))))
;; Reports memory allocations and potential leaks

;; Call stack inspection
(define (deep-function)
  (display-call-stack)
  (display-environment-variables))

;; Assertion debugging
(assert (> x 0) "x must be positive" x)
(debug-assert (sorted? lst) "list must be sorted" lst)

;; Interactive debugging
(debugger-break-on-error #t)
;; Drops into REPL when error occurs

;; Performance monitoring
(with-performance-monitor
  (lambda ()
    (large-computation))
  (lambda (stats)
    (display "Execution time: ") (display (stats 'time)) (newline)
    (display "Memory used: ") (display (stats 'memory)) (newline)))
\end{lambdustcode}
\end{example}

The standard library provides comprehensive functionality while maintaining R7RS compatibility and enabling advanced programming techniques through gradual adoption of Lambdust's extended features.